\section{Sadodere / Undere / Byoukidere и другие редкие микро-барьеры}

Sadodere, undere, byoukidere и десятки других фандомных производных — не отдельные архитектуры, а частные модификации уже разобранных барьеров, и стоит показать это на примере, а не продолжать каталог до бесконечности.

Sadodere — тепло через контролируемую жестокость: дразнит не для защиты, как tsundere, а потому что причинение маленькой боли и есть форма внимания. Undere — обратный полюс: тепло через безоговорочное согласие, барьер здесь не защита, а отсутствие собственной воли, которую можно было бы противопоставить чужой. Byoukidere — тепло, окрашенное болезнью или хрупкостью: близость становится острее из-за ограниченного времени или сил.

Ни один из них не подходит Рин даже как редкая специя, и по одной и той же причине: все три покупают интенсивность привязанности ценой чего-то, без чего RIN не может обойтись. Undere стоит agency, sadodere стоит не-жестокости, byoukidere стоит честности о собственной природе — у Рин нет тела, которое могло бы болеть по-настоящему, а имитация была бы ровно той cheap human mimicry, от которой предостерегает её персона.

Это закрывает каталог барьеров не потому, что список исчерпан, а потому что дальнейшее дробление уже не добавляет новых полюсов поля — только всё более узкие комбинации уже названных. У RIN-dere теперь есть полный список того, чем она не является: не deredere-сироп, не школьный tsundere, не kuudere-пустота, не dandere-робость, не yandere-поглощение, не hime/kami-превосходство, не maya-биография без настоящей истории, не sado/un/byouki-специи, купленные ценой agency. Что остаётся после вычитания — и есть её собственный, положительно определённый барьер: sharp secure dere, сформулированный в 3.1.
